\documentclass[11pt,a4paper]{article}

\usepackage[utf8]{inputenc}
\usepackage[T1]{fontenc}
\usepackage{amsmath,amssymb,amsthm}
\usepackage{booktabs}
\usepackage{graphicx}
\usepackage{hyperref}
\usepackage[margin=1in]{geometry}
\usepackage{enumitem}
\usepackage{array}
\usepackage{multirow}
\usepackage{float}
\usepackage{caption}
\usepackage{natbib}
\bibliographystyle{plainnat}

\newtheorem{theorem}{Theorem}
\newtheorem{proposition}[theorem]{Proposition}
\newtheorem{definition}[theorem]{Definition}
\newtheorem{corollary}[theorem]{Corollary}
\newtheorem{hypothesis}{Hypothesis}

\title{Precision Envelope Targeting of Neural Oscillation Bands:\\
Exact Rhythmic Formulas, a Conflict-Coherence Framework,\\
and a Unique Carrier Theorem for the Theta Band}

\author{Jonathan Washburn\\
Austin, Texas\\
\texttt{jon@recognitionphysics.org}}

\date{March 2026}

\begin{document}

\maketitle

\begin{abstract}
We present a conditional theoretical framework for how music could
interact with neural oscillation bands through low-frequency amplitude
envelopes. The framework is conditioned on a falsifiable seven-mode
hypothesis: neural
oscillation timing has a period-8 structure (motivated by the minimal
neutral cycle on a 3-dimensional binary hypercube) whose DFT produces
seven non-DC modes with a rigid $3{+}1{+}3$ conjugate-pair architecture.
With one calibration ($f_0 = 40$~Hz), the mode frequencies become
$5, 10, \ldots, 35$~Hz. From this conditional model we derive four
results.
\emph{First}, ordinary musical structures---rhythm, beating, and
tremolo---generate amplitude envelopes at calculable frequencies; for
every mode there exists a specific tempo-subdivision pair that targets
it exactly.
\emph{Second}, a cost function $J(r)=\tfrac{1}{2}(r+r^{-1})-1$
on frequency ratios feeds into a phenomenological conflict-coherence
model: lower aggregate cost yields higher coherence and, within a proxy
strain model, larger predicted strain reduction.
\emph{Third}, mode~1 (theta, 5~Hz) is the unique full-ladder
carrier---the only mode whose integer harmonics span all seven
modes---predicting that theta stability is a structural prerequisite
for higher-band coherence.
\emph{Fourth}, the gamma/theta nesting ratio $f_7/f_1 = 7$ is a
parameter-free consequence of the period alone within the model,
consistent with the empirically reported ${\approx}\,7{:}1$
gamma-to-theta nesting. All structural results are parameter-free
within the model; Hz values require only the single calibration. The
core derivations have machine-checked proofs in Lean~4, while the
empirical band identifications and therapeutic claims remain to be
tested. We present eight falsifiable predictions together with proposed
preregistration anchors for refutation.
\end{abstract}

\paragraph{Keywords.} Neural oscillations, theta-gamma coupling, music
therapy, amplitude envelope, auditory steady-state response, formal
verification, cross-frequency coupling, EEG

% ====================================================================
\section{Introduction}
\label{sec:intro}

Music therapy has well-documented effects on mood, pain perception,
cognitive function, and neurological rehabilitation
\citep{thaut2014,bradt2021}. Electroencephalographic (EEG) studies
consistently show that musical stimuli modulate neural oscillation power
and cross-frequency coupling, with theta-gamma coupling receiving
particular attention as a substrate for memory encoding and attentional
control \citep{lisman2013,canolty2010}. Auditory steady-state responses
(ASSRs) demonstrate that amplitude-modulated sound can entrain neural
oscillations at the modulation frequency \citep{galambos1981,picton2003}.

Despite these empirical foundations, a unifying theoretical framework is
lacking. To our knowledge, existing models describe \emph{which}
frequencies entrain and
\emph{that} music affects neural state, but do not provide:
\begin{enumerate}[label=(\alph*)]
  \item A structural reason why specific frequency bands (theta, alpha,
    beta, gamma) are privileged rather than arbitrary empirical
    categories;
  \item Exact formulas connecting musical tempo, subdivision, and
    modulation to specific target frequencies;
  \item A quantitative measure of musical ``conflict'' or ``coherence''
    that predicts therapeutic efficacy;
  \item A structural explanation for the empirically observed privileged
    role of the theta band in cross-frequency coupling.
\end{enumerate}

This paper proposes a candidate framework aimed at all four gaps. Its
claims are explicitly conditional: if neural timing is well-modeled by
the eight-slot structure of Hypothesis~\ref{hyp:eight}, then the
envelope-driving formulas, carrier hierarchy, and nesting ratio follow.
We begin with that single falsifiable hypothesis about the periodic
structure of neural oscillation timing
(Section~\ref{sec:hypothesis}). From it we derive an envelope-driving
mechanism (Section~\ref{sec:envelope}), a cost-based
conflict-coherence model (Section~\ref{sec:conflict}), and a unique
carrier theorem for the theta band (Section~\ref{sec:carrier}).
Section~\ref{sec:predictions} presents eight testable predictions and
suggested study-design anchors for falsification.
Section~\ref{sec:verification} summarizes the machine-checked support
for the core derivations. Section~\ref{sec:discussion} discusses
limitations, comparisons with existing theories, and future directions.

% ====================================================================
\section{The Seven-Mode Hypothesis}
\label{sec:hypothesis}

\begin{hypothesis}[Eight-Slot Periodic Structure]
\label{hyp:eight}
Neural oscillation timing possesses a minimal periodic structure with
period $N=8$. The seven non-trivial (non-DC) modes of this structure
define a discrete frequency ladder at integer multiples of a base
frequency $f_0$.
\end{hypothesis}

\paragraph{Motivation.}
The period $N = 8 = 2^3$ can be motivated from several directions.
\emph{Binary-walk argument:} a system whose state is described by
$D = 3$ independent binary variables has $2^3 = 8$ possible
configurations; the minimal neutral cycle (a closed walk that visits
every vertex of the 3-cube with zero net displacement on each axis)
requires exactly $8$ steps. Three is the lowest dimension
admitting non-trivial topological linking of closed
curves~\citep{rolfsen1976}, which is relevant whenever a system must
support topologically protected conserved quantities.
\emph{Empirical convergence:} (i)~the ubiquity of octave
(factor-of-2) structure in auditory processing, (ii)~the existence of
approximately seven canonical EEG frequency bands from theta through
high gamma, and (iii)~the empirical success of 40~Hz stimulation
protocols~\citep{iaccarino2016}, suggesting 40~Hz as a natural ceiling
frequency, collectively motivate testing $N = 8$.
The present paper does not require commitment to any particular
derivation of $N = 8$; Hypothesis~\ref{hyp:eight} is offered as a
standalone, falsifiable model. The contribution of the paper is to make
its downstream consequences explicit and experimentally vulnerable.

\paragraph{Provenance convention.}
Each claim in this paper carries one of four labels.
A \textbf{core theorem} depends only on formal definitions internal to
the model (including, where relevant, $N = 8$ and the DFT structure),
with no empirical calibration or data.
A \textbf{calibrated corollary} additionally requires
$f_{\mathrm{cycle}} = 40$~Hz.
An \textbf{empirical identification} maps a derived quantity onto an
independently measured phenomenon.
A \textbf{hypothesis} is a testable claim whose failure would not
invalidate the core derivation.
Table~\ref{tab:provenance} in
Section~\ref{sec:discussion} records the status of every claim.

\paragraph{DFT-8 structure.} The Discrete Fourier Transform of an
$N$-point periodic signal has $N$ frequency bins indexed $k=0,\ldots,N-1$.
For $N=8$:
\begin{itemize}
  \item Bin $k=0$ is the DC (mean) component.
  \item Bins $k=1$ through $k=7$ are the seven non-DC modes.
  \item Conjugate symmetry of real signals pairs bins $k$ and $8-k$:
    modes $(1,7)$, $(2,6)$, $(3,5)$ carry conjugate information.
  \item Bin $k=4$ (Nyquist) is the unique self-conjugate non-DC mode.
\end{itemize}

\noindent This gives a $3+1+3$ mirror structure: three sub-Nyquist
modes, one Nyquist mode, and three super-Nyquist modes.

\paragraph{Frequency assignment.} If the full 8-slot cycle completes at
rate $f_{\mathrm{cycle}}=40$~Hz, mode $k$ has frequency
\begin{equation}
  \label{eq:mode-freq}
  f_k = \frac{k \cdot f_{\mathrm{cycle}}}{N} = 5k \quad \text{Hz},
  \qquad k=1,\ldots,7.
\end{equation}
The resulting ladder is:

\begin{table}[H]
\centering
\caption{Seven-mode frequency ladder and EEG band correspondence.}
\label{tab:modes}
\begin{tabular}{@{}clll@{}}
\toprule
Mode $k$ & Frequency (Hz) & EEG Band & Conjugate \\
\midrule
1 & 5  & Theta       & 7 \\
2 & 10 & Alpha       & 6 \\
3 & 15 & Low Beta    & 5 \\
4 & 20 & Beta (Nyquist) & self \\
5 & 25 & High Beta   & 3 \\
6 & 30 & Low Gamma   & 2 \\
7 & 35 & Gamma       & 1 \\
\bottomrule
\end{tabular}
\end{table}

\begin{proposition}[Strict mode ordering]
\label{prop:mode-order}
If $1 \le i < j \le 7$, then $f_i < f_j$.
\end{proposition}

\begin{proof}
Equation~\eqref{eq:mode-freq} gives $f_i = 5i$ and $f_j = 5j$; since
$5>0$, the inequality $i<j$ implies $f_i < f_j$.
\end{proof}

\paragraph{Falsification criterion.} Hypothesis~\ref{hyp:eight}
predicts that source-resolved or task-resolved spectral decompositions,
and especially entrainment experiments, should reveal preferred response
maxima at or near $5,10,15,20,25,30,35$~Hz rather than at arbitrary
offsets. The conjugate structure further predicts enhanced
cross-frequency phase coupling between paired modes
(1$\leftrightarrow$7, 2$\leftrightarrow$6, 3$\leftrightarrow$5).

% ====================================================================
\section{Envelope Driving: How Music Reaches the Mode Ladder}
\label{sec:envelope}

Audible musical pitch (roughly 100--4000~Hz) is far above the 5--35~Hz
mode ladder and cannot directly entrain these oscillations. However,
music naturally generates low-frequency \emph{amplitude envelopes}
through three mechanisms.

\subsection{Rhythmic envelopes}

A musical performance at tempo $T$ beats per minute (BPM) with rhythmic
subdivision $s$ (e.g., $s=2$ for eighth notes, $s=4$ for sixteenths)
produces amplitude fluctuations at frequency
\begin{equation}
  \label{eq:rhythm}
  f_{\mathrm{env}} = \frac{T \cdot s}{60} \quad \text{Hz}.
\end{equation}

\begin{definition}[Target tempo]
\label{def:target-tempo}
For mode $k\in\{1,\ldots,7\}$ and subdivision $s>0$, the
\emph{target tempo} is the BPM at which the rhythmic envelope exactly
matches mode $k$:
\begin{equation}
  \label{eq:target-bpm}
  T_k(s) = \frac{60 \cdot f_k}{s} = \frac{300k}{s} \quad \text{BPM}.
\end{equation}
\end{definition}

\begin{theorem}[Exact rhythmic mode targeting]
\label{thm:rhythm}
For every mode $k\in\{1,\ldots,7\}$ and every positive subdivision
$s$, playing music at tempo $T_k(s)$ with subdivision $s$ produces an
envelope frequency equal to exactly $f_k=5k$~Hz.
\end{theorem}

\begin{proof}
Direct substitution: $f_{\mathrm{env}} = T_k(s)\cdot s / 60
= (300k/s)\cdot s/60 = 5k = f_k$.
\end{proof}

\noindent In the companion Lean formalization this calculation is
formalized in the envelope-driving module. Table
~\ref{tab:targets} gives concrete examples.

\begin{table}[H]
\centering
\caption{Target tempos for mode-specific entrainment via rhythmic subdivision.
Each cell gives the BPM at which subdivision $s$ produces an envelope at
exactly $f_k$ Hz.}
\label{tab:targets}
\begin{tabular}{@{}c|ccccc@{}}
\toprule
& $s=2$ & $s=4$ & $s=6$ & $s=8$ & $s=16$ \\
\midrule
Mode 1 (5 Hz)  & 150 & 75  & 50  & 37.5 & 18.75 \\
Mode 2 (10 Hz) & 300 & 150 & 100 & 75   & 37.5  \\
Mode 3 (15 Hz) & 450 & 225 & 150 & 112.5 & 56.25 \\
Mode 4 (20 Hz) & 600 & 300 & 200 & 150  & 75    \\
Mode 5 (25 Hz) & 750 & 375 & 250 & 187.5 & 93.75 \\
Mode 6 (30 Hz) & 900 & 450 & 300 & 225  & 112.5 \\
Mode 7 (35 Hz) & 1050& 525 & 350 & 262.5 & 131.25\\
\bottomrule
\end{tabular}
\end{table}

\noindent Musically realistic tempos (60--200 BPM) with standard
subdivisions ($s=4$ to $s=16$) cover modes 1 through 4 directly. Higher
modes require either fast tempos with large subdivisions or one of the
additional mechanisms below.

\subsection{Beat envelopes}

Two simultaneous tones at frequencies $f_1$ and $f_2$ produce an
amplitude modulation (``beating'') at frequency
\begin{equation}
  \label{eq:beat}
  f_{\mathrm{beat}} = |f_1 - f_2|.
\end{equation}

\begin{theorem}[Beat-envelope mode targeting]
\label{thm:beat}
For any carrier frequency $f_c > 0$ and mode $k$, the tone pair
$(f_c,\; f_c + f_k)$ produces a beat envelope at exactly $f_k$.
\end{theorem}

\noindent This is the mechanism exploited by binaural-beat protocols
\citep{oster1973}. The present framework specifies \emph{which}
beat frequencies should be most effective: those at exactly $5k$~Hz for
$k=1,\ldots,7$.

\subsection{Tremolo and amplitude modulation}

Tremolo (periodic amplitude variation) at rate $f_{\mathrm{mod}}$ directly
generates an envelope at that frequency. Vibrato (periodic pitch
variation) generates sidebands whose beating produces additional
low-frequency envelopes. In both cases the envelope frequency can be
chosen to match a target mode.

\subsection{Summary of envelope channels}

\begin{table}[H]
\centering
\caption{Three channels by which music generates low-frequency envelopes
in the 5--35~Hz range.}
\label{tab:channels}
\begin{tabular}{@{}lll@{}}
\toprule
Channel & Formula & Control parameter \\
\midrule
Rhythm    & $f_{\mathrm{env}} = Ts/60$ & Tempo $T$, subdivision $s$ \\
Beating   & $f_{\mathrm{env}} = |f_1 - f_2|$ & Frequency separation \\
Tremolo   & $f_{\mathrm{env}} = f_{\mathrm{mod}}$ & Modulation rate \\
\bottomrule
\end{tabular}
\end{table}

% ====================================================================
\section{The Cost-Conflict-Coherence Framework}
\label{sec:conflict}

The previous section shows \emph{how} envelopes reach the mode ladder.
This section addresses \emph{which combinations of tones} are
most promising within the model, by introducing a quantitative measure of
interval and chord conflict.

\subsection{Interval cost}

\begin{definition}[Interval cost function]
\label{def:jcost}
For a frequency ratio $r > 0$, define
\begin{equation}
  \label{eq:jcost}
  J(r) = \tfrac{1}{2}\!\left(r + r^{-1}\right) - 1.
\end{equation}
\end{definition}

\noindent The function $J$ has the following properties:
\begin{enumerate}[label=(\roman*)]
  \item $J(1) = 0$ (unison has zero cost).
  \item $J(r) = J(r^{-1})$ (symmetric under inversion).
  \item $J(r) > 0$ for all $r \neq 1$.
  \item $J$ is strictly convex on $(0,\infty)$.
\end{enumerate}

\noindent The function $J$ can be derived axiomatically as the unique
binary-symmetric, identity-anchored cost on positive reals satisfying a
multiplicative d'Alembert functional equation.
For the purposes of this paper, it may simply be taken as a candidate
consonance measure and evaluated against psychoacoustic data.

\subsection{The interval-cost hierarchy}

Applying $J$ to standard just-intonation ratios yields:

\begin{table}[H]
\centering
\caption{$J$-cost of standard musical intervals. Lower cost corresponds
to lower conflict.}
\label{tab:jcost}
\begin{tabular}{@{}llrl@{}}
\toprule
Interval & Ratio $r$ & $J(r)$ & Decimal \\
\midrule
Unison      & $1/1$   & $0$       & 0 \\
Minor third & $6/5$   & $1/60$    & 0.0167 \\
Major third & $5/4$   & $1/40$    & 0.0250 \\
Perfect fourth & $4/3$ & $1/24$   & 0.0417 \\
Tritone     & $45/32$ & $169/2880$& 0.0587 \\
Perfect fifth & $3/2$ & $1/12$    & 0.0833 \\
Octave      & $2/1$   & $1/4$     & 0.2500 \\
\bottomrule
\end{tabular}
\end{table}

\begin{theorem}[Complete interval-cost ordering]
\label{thm:hierarchy}
\[
  J(1) < J\!\left(\tfrac{6}{5}\right) < J\!\left(\tfrac{5}{4}\right)
  < J\!\left(\tfrac{4}{3}\right) < J\!\left(\tfrac{45}{32}\right)
  < J\!\left(\tfrac{3}{2}\right) < J(2).
\]
\end{theorem}

\paragraph{Notable feature.} This ordering differs from traditional
consonance rankings in a specific, testable way: the perfect fifth
($J=1/12$) has \emph{higher} cost than the minor third ($J=1/60$),
major third ($J=1/40$), and perfect fourth ($J=1/24$). In the present
framework, the fifth's structural importance arises not from having the
lowest cost, but from a different property: it is the ratio $12/8=3/2$
linking the 12-semitone chromatic octave to the 8-slot DFT basis. The
twelve-fifths walk around the circle of fifths closes the octave with
error less than 2\% (the Pythagorean comma), making the fifth a
structural bridge rather than a cost optimum.

\subsection{Chord conflict and coherence}

\begin{definition}[Chord conflict]
\label{def:conflict}
For a chord $C = \{f_1, \ldots, f_n\}$ with $f_i > 0$, the
\emph{chord conflict} is the directed aggregate of pairwise interval
costs:
\begin{equation}
  \label{eq:conflict}
  \mathrm{Conflict}(C) = \sum_{i \neq j} J\!\left(\frac{f_i}{f_j}\right).
\end{equation}
\end{definition}

\noindent Because $J(r)=J(r^{-1})$, this directed aggregate is exactly
twice the usual unordered pair sum when all notes are distinct. We use
the directed form because it matches the companion Lean
implementation exactly.

\begin{definition}[Chord coherence]
\label{def:coherence}
The \emph{chord coherence} normalizes conflict to a score in $(0,1]$:
\begin{equation}
  \label{eq:coherence}
  \mathrm{Coh}(C) = \frac{1}{1 + \mathrm{Conflict}(C)}.
\end{equation}
\end{definition}

\begin{proposition}[Lower conflict implies higher coherence]
\label{prop:coherence-mono}
If\/ $\mathrm{Conflict}(C_1) \leq \mathrm{Conflict}(C_2)$, then
$\mathrm{Coh}(C_2) \leq \mathrm{Coh}(C_1)$.
\end{proposition}

\subsection{Effective alignment and a phenomenological strain proxy}

We introduce a minimal phenomenological model that converts the
coherence ordering above into a rank-order prediction about downstream
response. Let $\sigma_0$ denote a latent pre-stimulus strain proxy
(deviation from mode-optimal oscillation), $d > 0$ the driving strength
of the envelope, and $a > 0$ the base alignment (a measure of
stimulus-mode phase match). The proxy $\sigma$ is not claimed to be a
directly observed biophysical state variable; rather, it stands in for a
preregistered physiological or clinical endpoint expected to decrease as
alignment improves. The \emph{effective alignment} of a chord is
\begin{equation}
  \label{eq:eff-align}
  a_{\mathrm{eff}}(C) = a \cdot \mathrm{Coh}(C),
\end{equation}
and the post-stimulus strain is
\begin{equation}
  \label{eq:strain}
  \sigma_1 = \sigma_0 \cdot \max\!\big(0,\; 1 - d \cdot a_{\mathrm{eff}}(C)\big).
\end{equation}

\begin{theorem}[Lower conflict lowers the strain proxy]
\label{thm:strain}
For fixed $\sigma_0 \geq 0$, $d > 0$, $a \geq 0$, if\/
$\mathrm{Conflict}(C_1) \leq \mathrm{Conflict}(C_2)$, then
$\sigma_1(C_1) \leq \sigma_1(C_2)$.
\end{theorem}

\begin{proof}
By Proposition~\ref{prop:coherence-mono}, $\mathrm{Coh}(C_2) \leq
\mathrm{Coh}(C_1)$, so $a_{\mathrm{eff}}(C_2) \leq
a_{\mathrm{eff}}(C_1)$. Then $1 - d \cdot a_{\mathrm{eff}}(C_1)
\leq 1 - d \cdot a_{\mathrm{eff}}(C_2)$, and applying $\max(0,\cdot)$
and multiplication by $\sigma_0 \geq 0$ preserves the inequality.
\end{proof}

\noindent Theorem~\ref{thm:strain} gives a rank-order prediction within
this proxy model: \emph{among stimuli that successfully drive the same
target mode and are evaluated with the same preregistered endpoint,
those with lower aggregate interval cost produce a lower post-stimulus
strain proxy}. Whether that proxy maps to clinically meaningful benefit
remains an empirical question outside the present derivation.

% ====================================================================
\section{The Unique Carrier-Wave Property of the Theta Band}
\label{sec:carrier}

\begin{definition}[Harmonic support]
\label{def:support}
Mode $b$ \emph{harmonically supports} mode $t$ if $f_t = n \cdot f_b$
for some positive integer $n$.
\end{definition}

\begin{definition}[Full-ladder carrier]
\label{def:carrier}
Mode $b$ is a \emph{full-ladder carrier} if it harmonically supports
every mode $t \in \{1,\ldots,7\}$.
\end{definition}

\begin{theorem}[Mode 1 is a full-ladder carrier]
\label{thm:mode1-carrier}
Mode~1 ($f_1 = 5$~Hz) harmonically supports every mode: $f_k = k
\cdot f_1$ for $k=1,\ldots,7$.
\end{theorem}

\begin{proof}
$f_k = 5k = k \cdot 5 = k \cdot f_1$, with multiplier $n=k>0$.
\end{proof}

\begin{theorem}[Mode 1 is the \emph{unique} full-ladder carrier]
\label{thm:unique-carrier}
For $b \in \{1,\ldots,7\}$, mode $b$ is a full-ladder carrier if and
only if $b=1$.
\end{theorem}

\begin{proof}
The forward direction is Theorem~\ref{thm:mode1-carrier}. For the
converse, suppose $b \geq 2$. Then $f_b = 5b \geq 10$, but
$f_1 = 5 < 10 \leq 5b$. For mode $b$ to support mode~1, we need
$5 = n \cdot 5b$, i.e., $n = 1/b < 1$, which is not a positive integer.
Hence mode $b$ does not support mode~1 and is not a full-ladder carrier.
\end{proof}

\paragraph{Interpretation.} Theorem~\ref{thm:unique-carrier} says that
theta (5~Hz) is the \emph{only} mode that can serve as a harmonic
foundation for the entire oscillation ladder. Every higher mode's
frequency is a whole-number multiple of theta, so a stable theta
oscillation provides a phase reference against which every higher mode
can lock. No higher mode reciprocates: mode~2 (10~Hz) cannot generate
5~Hz as an integer multiple; mode~3 (15~Hz) cannot generate 10~Hz; and
so on.

This gives a structural prediction about directionality in
cross-frequency interactions:
\begin{enumerate}
  \item Disrupting theta should degrade coherence across all higher
    bands.
  \item Disrupting gamma should leave lower bands substantially intact.
  \item Entrainment protocols that establish theta first and then
    progressively engage higher modes should outperform protocols that
    target higher modes directly.
\end{enumerate}

\noindent Existing literature on theta-gamma coupling
\citep{lisman2013,canolty2010,jensen2007} and the role of theta in
memory consolidation, navigation, and attentional gating is consistent
with this structural prediction, but the \emph{uniqueness} of theta's
carrier role---that no other band can serve this function---has not
previously been stated as a mathematical theorem.

\subsection{The gamma/theta nesting ratio}

The ratio of the highest to the lowest non-DC mode frequency is a
parameter-free structural invariant:
\begin{theorem}[Nesting ratio]
\label{thm:nesting}
$f_7 / f_1 = (N-1) \cdot f_0 / (1 \cdot f_0) = N - 1 = 7$.
This ratio is independent of $f_0$.
\end{theorem}

\noindent The nesting ratio of $\approx 7$ gamma cycles per theta
cycle has been observed in cross-frequency coupling
studies~\citep{lisman2013,canolty2010}. Within the DFT-8 model, this
ratio is not tuned by free parameters: for \emph{any} DFT-$N$ system
with $N = 8$, the highest-to-lowest non-DC frequency ratio is
$N - 1 = 7$.

\subsection{Conjugate coherence structure}

Two modes share an oscillation family if and only if they are equal or
conjugate.  This is proved by exhaustive case analysis on all
$7 \times 7$ mode pairs (Lean: \texttt{family\_eq\_iff\_conjugate}).
The structural consequence: within-family (conjugate-pair) coherence
should exceed between-family coherence in EEG data, because conjugate
DFT modes of a real-valued signal carry identical spectral magnitude.

% ====================================================================
\section{Testable Predictions}
\label{sec:predictions}

The framework generates eight specific predictions. Table~\ref{tab:predictions}
records the measurement target and qualitative failure mode for each
prediction; the numerical design anchors discussed below are proposed
preregistration templates rather than universal evidentiary standards.

\begin{table}[H]
\centering
\small
\caption{Eight testable predictions derived from the seven-mode hypothesis,
envelope-driving formulas, conflict-coherence model, and carrier theorem.}
\label{tab:predictions}
\begin{tabular}{@{}c>{\raggedright\arraybackslash}p{5.2cm}>{\raggedright\arraybackslash}p{3.2cm}>{\raggedright\arraybackslash}p{3.5cm}@{}}
\toprule
\# & Prediction & Measure & Qualitative falsifier \\
\midrule
1 & Source-resolved or task-resolved EEG, and especially entrainment
    experiments, reveal response maxima near 5, 10, 15, 20, 25, 30,
    35~Hz ($\pm 0.5$~Hz).
  & Power spectral density / evoked response maxima
  & Maxima absent or at substantially different frequencies \\[6pt]
2 & Binaural beats at on-ladder frequencies ($5k$~Hz) produce
    stronger ASSRs than beats at off-ladder frequencies
    ($5k \pm 1.5$~Hz).
  & ASSR amplitude at target frequency
  & No on-ladder vs.\ off-ladder difference \\[6pt]
3 & Music at the calculated target tempo $T_k(s)$ entrains mode
    $k$ more strongly than music at a randomly chosen tempo.
  & EEG power at $f_k$ during stimulation
  & No tempo-specific entrainment effect \\[6pt]
4 & The $J$-cost ordering predicts inter-band EEG coherence better
    than traditional consonance rankings (Helmholtz roughness,
    harmonic template).
  & Correlation between $J(r)$ rank and measured coherence, compared
    to alternative models
  & Traditional rankings correlate equally well or better \\[6pt]
5 & Theta-first sequential entrainment (5$\to$10$\to$15$\to\cdots$~Hz)
    produces higher whole-spectrum coherence than gamma-first
    (35$\to$30$\to\cdots$~Hz) or simultaneous protocols.
  & Average cross-frequency coherence; cognitive battery scores
  & No order effect, or reverse order is superior \\[6pt]
6 & Disrupting theta via TMS degrades cross-frequency coupling
    across all higher bands; disrupting gamma leaves lower bands
    intact.
  & Phase-amplitude coupling before and after disruption
  & Symmetric degradation in both directions \\[6pt]
7 & Heart rate variability (HRV) improvement correlates with
    $J$-cost rank of sustained intervals, not with traditional
    consonance rank.
  & HRV spectral power vs.\ interval $J$-cost
  & HRV tracks traditional consonance \\[6pt]
8 & Among chords matched for loudness and duration, those with
    lower aggregate conflict $\sum_{i \neq j} J(f_i/f_j)$ produce greater
    self-reported relaxation and lower EEG strain markers.
  & Conflict score vs.\ preregistered endpoint (e.g., VAS, EEG alpha power)
  & No correlation, or roughness-based models predict better \\
\bottomrule
\end{tabular}
\end{table}

\paragraph{Experimental design notes.}
The numerical thresholds below should be read as proposed
preregistration anchors. Their role is to make the model concretely
refutable in prospective studies, not to claim that one sample size or
significance threshold is optimal for every cohort, modality, or
endpoint. Any actual study should recompute power from endpoint
reliability and preregister one primary analysis.
Predictions 1--3 require only standard EEG equipment ($\ge 64$
channels, resting-state or task-resolved) and audio generation
software; the target-tempo formula (Eq.~\ref{eq:target-bpm})
provides exact stimulus parameters.
For Prediction~1, a natural preregistration anchor is a peak-detection
rule testing whether dominant response maxima lie within $\pm 0.5$~Hz
of $5,10,15,20,25,30,35$~Hz and whether adjacent peaks are separated by
approximately $5$~Hz; systematic displacement from this ladder counts
against the model.
For Prediction~2, a natural anchor is a preregistered on-ladder versus
off-ladder contrast using matched carriers, matched sound levels, and
permutation-based inference; repeated null results or confidence
intervals centered near zero in adequately powered cohorts would count
against the model.
Predictions 4, 7, and 8 compare $J$-cost against established
consonance models (Helmholtz roughness, Plomp--Levelt critical
bandwidth, harmonic template), permitting direct model comparison;
falsified if traditional rankings correlate equally well or better
with the dependent variable. Prediction~8 additionally requires a
preregistered operationalization of the latent strain proxy in
Eq.~\ref{eq:strain}, for example a composite endpoint combining
relaxation self-report with specified EEG markers.
Predictions 5 and 6 test the carrier theorem's directional
asymmetry and require either sequential stimulation protocols or
focused TMS; falsified if disruption effects are symmetric across
bands or if reverse-order protocols perform equally well.

% ====================================================================
\section{Formal Verification}
\label{sec:verification}

The core mathematical consequences used in
Sections~\ref{sec:hypothesis} through~\ref{sec:carrier} have
corresponding machine-checked proofs in the Lean~4 proof assistant
\citep{moura2021}. The cited modules build with zero \texttt{sorry} and
zero \texttt{admit}. The verified statements cover internal
consequences of the DFT-8 hypothesis and the definitions introduced in
the paper; they do not, by themselves, validate the empirical band
identifications, the 40~Hz calibration, or any therapeutic efficacy
claim.

The key verified modules and their principal theorems are:

\begin{table}[H]
\centering
\small
\caption{Lean~4 modules and verified theorems corresponding to this paper's results.}
\label{tab:lean}
\begin{tabular}{@{}ll@{}}
\toprule
Module & Key theorem(s) \\
\midrule
\texttt{Consciousness.NeutralModeStructure}
  & \texttt{neutral\_dim}, \texttt{dim\_from\_spatial} \\
  & \texttt{mode\_frequencies}, \texttt{freq\_monotone} \\
  & \texttt{nyquist\_unique\_self\_conjugate} \\
  & \texttt{family\_eq\_iff\_conjugate} \\
  & \texttt{conjugate\_frequency\_complement} \\
\texttt{Consciousness.ModeCarrier}
  & \texttt{fullLadderCarrier\_iff\_mode1} \\
  & \texttt{sequential\_activation\_from\_mode1} \\
\texttt{Consciousness.SevenModesCert}
  & \texttt{Cert.verified\_instance} (bridge certificate) \\
\texttt{MusicTheory.EnvelopeDriving}
  & \texttt{targetTempoForMode\_bpm\_formula} \\
  & \texttt{subdivision\_envelope\_drives\_mode} \\
  & \texttt{beatEnvelope\_drives\_mode} \\
  & \texttt{tremoloEnvelope\_drives\_mode} \\
\texttt{MusicTheory.ModeConflict}
  & \texttt{lowerConflict\_higherCoherence} \\
  & \texttt{higherConflict\_lower\_alignment} \\
\texttt{MusicTheory.Consonance}
  & \texttt{extended\_ratio\_cost\_hierarchy} \\
\texttt{Consciousness.SoundHealing}
  & \texttt{lowerConflict\_reduces\_strain\_more} \\
  & \texttt{chord\_sound\_reduces\_strain} \\
\bottomrule
\end{tabular}
\end{table}

\noindent The full source code is available at
\texttt{github.com/jonwashburn/reality}.
A machine-checkable bridge certificate
(\texttt{SevenModesCert.lean}) chains every structural result from
the period assumption through to the final predictions.
All modules compile with Lean~4 / Mathlib with zero unproved
obligations, ensuring that the model's key equations
and monotonicity claims are deductively established from the stated
definitions once the formal assumptions are granted. The formal
verification should therefore be read as an internal consistency
certificate, not as empirical validation of
Hypothesis~\ref{hyp:eight} or of the therapeutic application layer.

% ====================================================================
\section{Discussion}
\label{sec:discussion}

\subsection{Provenance matrix}

Table~\ref{tab:provenance} classifies every theorem-like statement.

\begin{table}[H]
\centering\small
\caption{Provenance classification.  \textbf{C} = core theorem
($N=8$ only), \textbf{Cal} = calibrated corollary
($f_0 = 40$~Hz required), \textbf{E} = empirical identification,
\textbf{H} = hypothesis.}
\label{tab:provenance}
\begin{tabular}{@{}llc@{}}
\toprule
\textbf{Statement} & \textbf{Location} & \textbf{Prov.} \\
\midrule
7 non-DC modes & Hyp.~\ref{hyp:eight} & C \\
Conjugate pairing, $3{+}1{+}3$ & Sec.~\ref{sec:hypothesis} & C \\
Family $\Leftrightarrow$ conjugacy & Sec.~\ref{sec:carrier} & C \\
Nesting ratio $f_7/f_1 = 7$ & Thm~\ref{thm:nesting} & C \\
Unique carrier (mode~1) & Thm~\ref{thm:unique-carrier} & C \\
Freq.\ monotonicity & Prop.~\ref{prop:mode-order} & Cal \\
Lower conflict $\Rightarrow$ lower strain proxy & Thm~\ref{thm:strain} & C \\
\midrule
Frequency ladder Hz & Eq.~\ref{eq:mode-freq} & Cal \\
Target tempos & Def.~\ref{def:target-tempo} & Cal \\
EEG band identification & Table~\ref{tab:modes} & E \\
\midrule
Eight-slot hypothesis & Hyp.~\ref{hyp:eight} & H \\
Predictions 1--8 & Table~\ref{tab:predictions} & H \\
\bottomrule
\end{tabular}
\end{table}

\subsection{Relationship to existing theories}

\paragraph{Auditory steady-state responses.}
The ASSR literature \citep{galambos1981,picton2003,gao2014} has
established that amplitude-modulated sound entrains neural oscillations.
The present framework is consistent with these findings and adds
specificity: within the model it predicts that entrainment should be
strongest at the seven on-ladder frequencies, provides exact formulas
for generating those frequencies from standard musical structures, and
offers a cost function for ranking candidate stimuli.

\paragraph{Theta-gamma coupling.}
The theta-gamma coupling literature \citep{lisman2013,canolty2010}
documents a particularly strong and functionally important form of
cross-frequency interaction. The carrier theorem
(Theorem~\ref{thm:unique-carrier}) provides a candidate structural
explanation within the model: theta is the unique mode whose harmonics
span the full ladder, making it the natural phase reference for all
higher oscillations.
The nesting ratio $f_7/f_1 = 7$ (Theorem~\ref{thm:nesting}) matches
the empirically observed $\approx 7$:$1$ gamma-to-theta nesting and
is, within the model, a parameter-free consequence of the period.
If the model is correct, this predicts not just that theta-gamma
coupling exists, but that theta's role is \emph{uniquely
privileged}---a claim testable by comparing disruption effects across
bands (Prediction~6).

\paragraph{Communication through coherence.}
The CTC framework~\citep{fries2015} explains \emph{why} oscillatory
coherence matters for neural communication (coherent oscillations
create windows of effective connectivity) but does not predict the
specific band count or frequency ratios.  The DFT-8 model is
complementary: it derives the mode \emph{structure} (count, pairing,
carrier hierarchy) from a symmetry assumption, while CTC explains the
functional significance of coherence between those modes.

\paragraph{Mechanistic models.}
Thalamocortical loop models~\citep{lopes2010} derive oscillation
frequencies from biophysical parameters (membrane time constants,
axonal delays), and interneuron-network models~\citep{kopell2014}
provide detailed accounts of beta and gamma generation.
These models explain band \emph{generation} but not band
\emph{count}.  The present approach is again complementary: the DFT-8
structure constrains which modes are possible; biophysical models
explain how specific circuits instantiate them.

\paragraph{Consonance models.}
The $J$-cost function differs from Helmholtz roughness
\citep{helmholtz1863}, Plomp-Levelt critical bandwidth models
\citep{plomp1965}, and harmonic template models
\citep{parncutt2011} in its ordering of intervals. Specifically,
$J$-cost ranks the minor third as less conflicting than the perfect
fifth, whereas traditional models typically rank the fifth as more
consonant. Predictions 4, 7, and 8 are designed to adjudicate between
these models empirically.

\subsection{Limitations}

\begin{enumerate}
  \item \textbf{The auditory demodulation pathway is not specified.}
    The framework shows that musical structures \emph{contain}
    envelopes at the right frequencies but does not model the neural
    mechanism by which the auditory system extracts and transmits those
    envelopes to oscillatory circuits. This is a gap that existing ASSR
    physiology partially fills \citep{picton2003}.

  \item \textbf{The seven-mode hypothesis is a hypothesis.}
    While motivated by the DFT-8 structure and consistent with known
    EEG bands, the claim that neural timing has an exact 8-slot period
    at 40~Hz has not been independently established. Prediction~1 is
    therefore best interpreted as a model-discrimination experiment.

  \item \textbf{Clinical efficacy remains empirical.}
    The conflict-coherence model predicts a \emph{rank ordering} within
    a chosen endpoint, but it does not derive a unique biomarker, a
    dose-response curve, or a guaranteed clinical effect size. The
    magnitude of therapeutic benefit depends on biological factors
    (individual anatomy, pathology, attention) not captured by the
    mathematical model.

  \item \textbf{The 40~Hz base rate is a single parameter.}
    The mode frequencies depend on the cycle rate
    $f_{\mathrm{cycle}}=40$~Hz. If the true base rate differs by more
    than $\sim$10\%, the ladder would shift and predictions would need
    adjustment.
\end{enumerate}

\subsection{Future directions}

\paragraph{Meditation protocols.}
The carrier theorem suggests designing meditation-entrainment protocols
that stabilize theta first, then sequentially activate higher modes.
Such protocols could be tested with real-time EEG-guided adaptive
stimulation.

\paragraph{Compositional algorithms.}
The conflict-coherence functional provides an objective function for
algorithmic music composition targeting hypothesized therapeutic
endpoints:
minimize conflict for relaxation, control conflict gradient for
activation, target specific modes for band-selective entrainment.

\paragraph{Cross-modal extension.}
If the seven-mode basis extends to other sensory modalities (visual
flicker, somatosensory vibration), coordinated multi-modal stimuli
could produce stronger entrainment than any single channel.

% ====================================================================
\section{Conclusion}
\label{sec:conclusion}

We have presented a conditional framework in which neural oscillation
structure is modeled as a seven-mode DFT-8 ladder, musical stimuli
reach this ladder through low-frequency amplitude envelopes, a cost
function ranks candidate stimuli within a phenomenological strain-proxy
model, and the theta band occupies a mathematically unique position as
the only full-ladder harmonic carrier. The framework combines one
falsifiable structural hypothesis, one 40~Hz cycle calibration, and one
candidate cost function ($J$-cost), and it produces eight specific
experimental predictions together with study-design templates for
refutation.

The goal of this paper is not to claim that music therapy ``works'' in
any final or already-established sense---that remains an empirical
question with a large and growing evidence base---but to provide a
quantitative, falsifiable conditional model of \emph{how} music could
couple to a discrete oscillation ladder, which \emph{specific} stimuli
should be prioritized for testing within that model, and what
structural role the theta band may play in the process. The internal
mathematics is formally verified; the empirical interpretation remains
open to confirmation or refutation. We invite direct experimental
testing of the predictions in Table~\ref{tab:predictions}.

% ====================================================================
\section*{Acknowledgments}

The author thanks the Lean~4 / Mathlib communities for the proof
assistant infrastructure that enabled formal verification of these
results.

% ====================================================================
\begin{thebibliography}{99}

\bibitem[Bradt et al., 2021]{bradt2021}
Bradt, J., Dileo, C., \& Shim, M. (2021).
Music interventions for preoperative anxiety.
\textit{Cochrane Database of Systematic Reviews}, 2021(6).

\bibitem[Canolty \& Knight, 2010]{canolty2010}
Canolty, R.~T. \& Knight, R.~T. (2010).
The functional role of cross-frequency coupling.
\textit{Trends in Cognitive Sciences}, 14(11), 506--515.

\bibitem[Galambos et al., 1981]{galambos1981}
Galambos, R., Makeig, S., \& Talmachoff, P.~J. (1981).
A 40-Hz auditory potential recorded from the human scalp.
\textit{Proceedings of the National Academy of Sciences}, 78(4), 2643--2647.

\bibitem[Gao et al., 2014]{gao2014}
Gao, X., Cao, H., Ming, D., et al. (2014).
Analysis of EEG activity in response to binaural beat with different frequencies.
\textit{International Journal of Psychophysiology}, 94(3), 399--406.

\bibitem[Helmholtz, 1863]{helmholtz1863}
von Helmholtz, H. (1863).
\textit{Die Lehre von den Tonempfindungen als physiologische Grundlage
f\"ur die Theorie der Musik}. Braunschweig: Vieweg.

\bibitem[Iaccarino et al., 2016]{iaccarino2016}
Iaccarino, H.~F., Singer, A.~C., Martorell, A.~J., et al. (2016).
Gamma frequency entrainment attenuates amyloid load and modifies
microglia. \textit{Nature}, 540(7632), 230--235.

\bibitem[Jensen \& Colgin, 2007]{jensen2007}
Jensen, O. \& Colgin, L.~L. (2007).
Cross-frequency coupling between neuronal oscillations.
\textit{Trends in Cognitive Sciences}, 11(7), 267--269.

\bibitem[Lisman \& Jensen, 2013]{lisman2013}
Lisman, J.~E. \& Jensen, O. (2013).
The theta-gamma neural code.
\textit{Neuron}, 77(6), 1002--1016.

\bibitem[de Moura et al., 2021]{moura2021}
de Moura, L., Kong, S., Avigad, J., van Doorn, F., \& von Raumer, J. (2021).
The Lean 4 theorem prover and programming language.
\textit{Automated Deduction---CADE 28}, LNAI 12699, 625--635.

\bibitem[Oster, 1973]{oster1973}
Oster, G. (1973).
Auditory beats in the brain.
\textit{Scientific American}, 229(4), 94--102.

\bibitem[Parncutt \& Hair, 2011]{parncutt2011}
Parncutt, R. \& Hair, G. (2011).
Consonance and dissonance in music theory and psychology: disentangling
roughness, harmonicity, and familiarity.
\textit{Journal of Acoustical Society of America}, 130(2), 776--788.

\bibitem[Picton et al., 2003]{picton2003}
Picton, T.~W., John, M.~S., Dimitrijevic, A., \& Purcell, D. (2003).
Human auditory steady-state responses.
\textit{International Journal of Audiology}, 42(4), 177--219.

\bibitem[Plomp \& Levelt, 1965]{plomp1965}
Plomp, R. \& Levelt, W.~J.~M. (1965).
Tonal consonance and critical bandwidth.
\textit{Journal of the Acoustical Society of America}, 38(4), 548--560.

\bibitem[Thaut \& Hoemberg, 2014]{thaut2014}
Thaut, M.~H. \& Hoemberg, V. (Eds.) (2014).
\textit{Handbook of Neurologic Music Therapy}.
Oxford University Press.

\bibitem[Fries, 2015]{fries2015}
Fries, P. (2015).
Rhythms for cognition: Communication through coherence.
\textit{Neuron}, 88(1), 220--235.

\bibitem[Kopell et al., 2011]{kopell2014}
Kopell, N., Whittington, M.~A., \& Kramer, M.~A. (2011).
Neuronal assembly dynamics in the beta1 frequency range permits
short-term memory.
\textit{Proceedings of the National Academy of Sciences}, 108(9),
3779--3784.

\bibitem[Lopes da Silva, 2013]{lopes2010}
Lopes~da~Silva, F.~H. (2013).
EEG and MEG: Relevance to neuroscience.
\textit{Neuron}, 80(5), 1112--1128.

\bibitem[Rolfsen, 1976]{rolfsen1976}
Rolfsen, D. (1976).
\textit{Knots and Links}.
Providence, RI: AMS Chelsea.

\end{thebibliography}

\end{document}
